The three reallocation tests use the 2015 Market Stability Reserve (MSR) reform as their identifying event. Decision (EU) 2015/1814 of October 2015 established a rules-based mechanism that removes allowances from circulation whenever the cumulative number of allowances in circulation exceeds 833 million tCO\textsubscript{2}, and reinjects allowances when it falls below 400 million \citep{Perino2018}. The reserve became operational in January 2019. The realized price path is in Figure~\ref{fig:eua_price_timeseries}: EUA prices rose from \EUR{5}--\EUR{8} per tCO\textsubscript{2} in 2015--16 to above \EUR{25} by 2018 and above \EUR{80} in 2022.

The MSR is plausibly exogenous to Belgian firm-level conditions for two reasons. First, the policy was the outcome of EU-level political negotiations in the years preceding 2015 and was not a response to Belgian conditions; the trigger threshold and the intake-rate parameters were fixed in the 2015 decision, before the price rise that the policy ultimately produced. Second, once the policy was enacted, the post-2018 EUA path is determined by the accumulated allowance surplus --- a deterministic function of past emissions and free allocation --- and the MSR's mechanical intake schedule. The path is therefore policy-driven rather than driven by short-run macroeconomic or industrial conditions \citep{DeJongheEtAl2021}.\footnote{The mechanical character of the MSR has been used as the identifying restriction in several recent papers on the post-2017 EU ETS price path; see also \citet{kanzig2023unequal} and \citet{BauerKaenzigRudebusch2026}.}

We use 2015 as the headline cutoff for the MSR tests. The within-NACE-4d analysis (\S\ref{sec:domestic_within_intensive}) reports three event-window cutoffs --- Phase~II onset in 2008, Phase~III onset in 2013, and the MSR in 2017 --- to bracket the MSR results against the pre-MSR phase transitions. The international margin (\S\ref{sec:international}) reports the share of regulated imports from EU and non-ETS sources from 2000 through 2022.
